\section{Resource Scheduling}
\label{sec:resource-scheduling}

Resource scheduling concerns the assignment of limited resources to tasks over time, subject to constraints on availability, task requirements, and temporal dependencies \parencite{pinedo2016scheduling}. \textcite{pinedo2016scheduling} defines the general scheduling problem as determining when and how to process a set of jobs using a set of machines, where each job has processing requirements and each machine has capacity limits. In transport operations, the ``machines'' are drivers and vehicles, and the ``jobs'' are assignments with fixed time windows. This structure appears across several related domains. Nurse scheduling matches staff to shifts based on qualifications and availability. Crew scheduling in airlines assigns pilots and cabin crew to flight legs under regulatory rest requirements. Driver scheduling in public transit pairs operators with runs across a service day. All of these problems share the same core structure: matching a pool of constrained resources to a set of time-bound tasks.

Ressursplanlegger adds a layer of complexity by assigning \emph{both} an employee and a vehicle to each assignment. This makes it a multi-resource scheduling problem. Where a single-resource problem (e.g., assigning only drivers) must search through $n \times m$ possible pairings of $n$ drivers and $m$ assignments, a two-resource problem must also select a vehicle for each pairing. The number of feasible combinations grows quickly. In the Ressursplanlegger domain, each assignment is a task with a fixed time window and a set of resource requirements: a driver with the right competencies and availability, and a vehicle of the correct type that is not already in use. The objective is to maximise assignment coverage while balancing several secondary goals, such as even workload distribution and respect for driver preferences.

The constraints in this problem fall into two categories. Hard constraints must be satisfied for a plan to be feasible. If any hard constraint is violated, the plan is invalid regardless of how well it performs on other measures. In Ressursplanlegger, the hard constraints are: an employee must hold all competencies required by the assignment (HC-01); the employee must be available and not already assigned (HC-02); the vehicle must be in active status (HC-03); the vehicle type must match the assignment's requirement (HC-04); and neither the employee (HC-05) nor the vehicle (HC-06) may be double-booked. Soft constraints, by contrast, represent preferences that the algorithm tries to satisfy but may violate at a cost. Each soft constraint carries a configurable weight that determines how much its violation affects the total solution score. The soft constraints include preferring an employee's designated vehicle (weight 10), balancing workload evenly (weight 20), minimising transitions between consecutive assignments (weight 5), scheduling high-priority assignments first (weight 15), and respecting stated employee preferences (weight 10). This distinction between hard and soft constraints follows the standard constraint satisfaction framework described by \textcite{rossi2006handbook}, where a feasible solution satisfies all hard constraints and an optimal solution additionally minimises the weighted sum of soft constraint violations.

Resource scheduling problems of this kind are NP-hard at realistic fleet sizes \parencite{pinedo2016scheduling}. An exact solver that guarantees an optimal solution must, in the worst case, explore an exponential number of combinations. For small instances this is tractable, but as the number of assignments, drivers, and vehicles grows, computation time becomes impractical. This motivates the use of heuristic and metaheuristic methods that trade optimality guarantees for speed. Ressursplanlegger addresses this with a multi-engine approach: three solvers, each suited to different problem sizes. A related theoretical area is the Vehicle Routing Problem (VRP), first formalised by \textcite{dantzig1959truck} and surveyed extensively by \textcite{toth2014vrp}. VRP and its time-windowed variant (VRPTW) share structural similarities with driver-vehicle assignment: both involve assigning resources to geographically distributed tasks under time and capacity constraints. However, the problems diverge in one important respect. VRP is primarily a sequencing problem: the algorithm decides the order in which a vehicle visits locations, and routing cost depends on travel distances between stops. In the Ressursplanlegger domain, assignment times and locations are fixed by the customer. The algorithm decides \emph{who} performs each assignment, not \emph{in which order}. Sequencing is not part of the problem.

The three solvers in Ressursplanlegger occupy different points on the speed-quality tradeoff. The greedy solver sorts assignments by priority and applies a first-fit heuristic, assigning each task to the first driver-vehicle pair that satisfies all hard constraints. It runs in $O(n \times m)$ time, where $n$ is the number of assignments and $m$ is the number of available resources, and produces a result almost instantly. It offers no optimality guarantee. The CP-SAT solver, based on Google OR-Tools, is a complete constraint programming solver that models the full set of hard and soft constraints and searches for the best feasible solution within a configurable time limit (default 60 seconds). It produces near-optimal results for instances of up to roughly 500 assignments. For larger instances (1\,000+ assignments), the system uses Timefold, a Java-based metaheuristic engine that combines tabu search, simulated annealing, and late-acceptance strategies to explore the solution space. Timefold does not guarantee optimality but scales better than CP-SAT on large problems. The theoretical basis for constraint programming as a solving paradigm is described by \textcite{rossi2006handbook}, who characterise it as a method for modelling and solving problems defined by variables, domains, and constraints. This multi-solver architecture allows the traffic coordinator to choose the appropriate solver for the problem at hand, trading computation time for solution quality depending on the day's complexity.
